% ------------------------------------------------------------
% CHAPTER 26
% ------------------------------------------------------------

\chapter{Lamphrodynamic Relaxation and Entropy Descent}

Delta–sigma modulation shapes error through feedback, routing quantization disturbances into high-frequency modes that do not interfere with the signal of interest. In earlier chapters, this process was understood in terms of noise transfer functions, geometric connections, and information-theoretic free-energy minimization. In this chapter, we reveal a deeper structure: delta–sigma modulators implement lamphrodynamic relaxation, a field-theoretic process in RSVP where entropy spreads outward, smoothing local disturbances and restoring coherence to the field. The modulator becomes a discretized lamphrodynamic engine whose internal dynamics are governed by entropy descent and curvature-damped flows. This interpretation completes the unification of delta–sigma theory with RSVP dynamics.

In the RSVP framework, the lamphrodynamic equation describes the evolution of the entropy field \(S\) coupled to the scalar field \(\Phi\) and vector field \(v\). In its simplest form, the entropy field satisfies a relaxation equation of the form
\[
\frac{\partial S}{\partial t} = \nabla \cdot ( D \nabla S ) - \kappa S + \sigma(t),
\]
where \(D\) is a diffusion coefficient, \(\kappa\) is a dissipation constant, and \(\sigma(t)\) represents localized entropy injections. In physical systems, this equation governs the spreading of disorder away from a perturbation, returning the field to a state of coherence. The similarity to delta–sigma modulation is immediate: quantization produces localized entropy injections, the loop filter forces the entropy into higher-order modes where it dissipates, and the state trajectory returns to coherence.

To connect this explicitly to delta–sigma modulation, consider the quantization error sequence \(e[n]\). At each quantization event, the scalar field \(x[n]\) undergoes a collapse that produces a discontinuity in the entropy field. This discontinuity is instantaneous and localized, analogous to the entropy injection term \(\sigma(t)\). The loop filter then acts as a smoothing operator, implementing a discrete convolution that approximates integration or higher-order integration. This smoothing corresponds to the diffusion term in the lamphrodynamic equation. Thus the loop filter becomes the mechanism through which entropy spreads away from the low-frequency region of interest, preventing local accumulation of disorder.

The delta–sigma loop equation,
\[
x[n+1] = x[n] + u[n] - y[n],
\]
can be interpreted as a discrete lamphrodynamic descent. The quantizer output \(y[n]\) introduces a jump in the entropy field because the collapse \(x[n] \mapsto Q(x[n])\) is discontinuous. The difference \(u[n] - y[n]\) represents the error that must be corrected by the smoothing process. Over time, the recurrence relation smooths the perturbation by integrating the error. This is exactly the behavior of lamphrodynamic relaxation, where perturbations decay under the action of a smoothing operator.

To formalize this connection, define the discrete entropy field \(S[n]\) as a function of the quantization error. For a one-bit quantizer,
\[
S[n] = | e[n] |,
\]
while for multi-bit quantizers, one may define entropy as the divergence between the continuous input and the quantized symbol in an information-geometric sense. The evolution of \(S[n]\) follows from the evolution of the error:
\[
S[n+1] = f(S[n], e[n], x[n]).
\]
Under mild assumptions of small error, the linearized dynamics of \(S\) approximate a discrete diffusion equation,
\[
S[n+1] - S[n] \approx -\nabla \cdot J_{S}(n),
\]
where \(J_{S}\) is an entropy flux proportional to the derivative of the loop filter. A first-order modulator produces a flux analogous to a first-order diffusion. Higher-order modulators produce higher-order diffusion or even anti-diffusion terms near the Nyquist frequency, which is why high-order modulators exhibit intricate spectral behavior. The entropic interpretation clarifies this: high-order modulation corresponds to extending the entropy-diffusion mechanism into higher spatial harmonics.

The most compelling aspect of lamphrodynamic relaxation in delta–sigma loops is the presence of an entropic Lyapunov functional. In RSVP, entropy descent is characterized by the monotonic decrease of a functional such as
\[
\mathcal{L}[S] = \int S^{2}(t) \, dt,
\]
whose gradient flow describes the natural dissipation of disorder. A delta–sigma modulator possesses a similar functional. Define
\[
\mathcal{V} = \sum_{n} S[n]^{2},
\]
which measures the squared magnitude of the error ripple. For stable modulators, \(\mathcal{V}\) decreases over time for any bounded input. This decreasing behavior matches the lamphrodynamic principle: the system moves toward coherence by dissipating localized entropy back into the high-frequency region.

The integrator in a first-order modulator corresponds to the simplest lamphrodynamic flow. Its impulse response spreads error over time, creating a smoothing effect akin to entropy diffusion. Higher-order modulators correspond to repeated applications of this smoothing operator. In continuous-time modulators, the loop filter becomes a differential operator that directly implements a lamphrodynamic PDE on the error field:
\[
\frac{d^{L} x}{d t^{L}} = u(t) - Q(x(t)).
\]
Viewed through lamphrodynamics, the right-hand side acts as an entropy source, while the left-hand side defines a PDE-driven relaxation flow. In the limit of high oversampling, the solution to this PDE approximates a low-curvature trajectory that filters out high-frequency entropy.

This theory also explains overload behavior in a new light. When the input amplitude exceeds the stable operating range, quantization induces large spikes in the entropy field. The lamphrodynamic flow may not dissipate such spikes quickly enough to preserve coherence. The system becomes trapped in a high-entropy basin, leading to limit cycles, tonal behavior, or instability. Overload recovery corresponds to the system re-entering the lamphrodynamic basin in which entropy can be dissipated effectively. The speed of recovery depends on the shape of the free-energy landscape and the operator norm of the smoothing filter.

Stability analysis becomes simpler in the lamphrodynamic picture. Instead of analyzing pole locations or spectral radius, one studies the monotonicity of the entropy functional under the loop dynamics. If the lamphrodynamic flow always decreases the entropy functional except at equilibrium, the system is globally asymptotically stable. This condition yields a geometric stability criterion:
\[
\frac{d}{dt} \mathcal{V}(t) \leq 0.
\]
This inequality is a discrete version of the lamphrodynamic Lyapunov condition, ensuring that every perturbation decays rather than grows. Delta–sigma stability becomes a form of lamphrodynamic entropic contractivity.

The unified picture that emerges is the following: delta–sigma modulators are lamphrodynamic relaxation engines that repeatedly convert local disorder into harmless high-frequency ripples. They operate by injecting entropy through quantizer collapse, smoothing it through integrator dynamics, and dissipating it through spectral shaping. Their stability and resolution arise from the geometry of entropy descent. Their accuracy arises from the smoothness of the scalar field. Their nonlinear dynamics arise from the interplay between lamphrodynamic diffusion and quantization-driven spikes.

The connection between lamphrodynamics and delta–sigma modulation is not metaphorical; it is mathematical. The modulator is a discretized entropy-dissipation PDE solver. It is therefore unsurprising that its behavior mirrors the RSVP dynamics of physical fields. This understanding prepares us to examine the deeper recursive structures underlying delta–sigma loops, which naturally bring us to the CLIO framework of recursive coherence and multi-level inference.


